\documentclass[11pt]{article}
\input{style-pc}

\begin{document}

\entetesujet{Sujet bac 2013 -- Physique-Chimie -- Série D}

% ============================ CHIMIE ============================
\matiere{CHIMIE}{8 points}

\exo{1}{(Spectre de l'atome d'hydrogène)}
Les niveaux d'énergie quantifiés de l'atome d'hydrogène sont donnés par :
\[ E_n = -\frac{13{,}6}{n^2}\ (\text{eV}) \quad n \text{ étant un nombre entier supérieur ou égal à } 1. \]

\begin{enumerate}
  \item \begin{enumerate}[label=\alph*.]
    \item Calculer les énergies des niveaux correspondant à $n = 1$ ; $n = 2$ ; $n = 3$ ; $n = 4$ ; $n = 5$ et $n = \infty$.
    \item Représenter dans le diagramme d'énergie les cinq premiers niveaux d'énergie ainsi que le niveau correspondant à l'atome ionisé.
  \end{enumerate}
  \item Le spectre de l'atome d'hydrogène contient les radiations de fréquences $N_a = 6{,}16\cdot10^{14}$ Hz et $N_b = 6{,}91\cdot10^{14}$ Hz. Sachant que ces deux radiations aboutissent au niveau $n = 2$, déterminer les numéros $a$ et $b$ des niveaux initiaux.
\end{enumerate}
\noindent On donne : $h = 6{,}62\cdot10^{-34}$ J$\cdot$s ; $1$ eV $= 1{,}6\cdot10^{-19}$ J.

\exo{2}{(Réaction acide-base)}
On dissout une quantité de soude de masse $m = 4{,}0$ g dans un volume $V_1 = 500$ mL de solution d'acide chlorhydrique de concentration molaire volumique $C_1 = 10^{-1}$ mol/L.

\begin{enumerate}
  \item \begin{enumerate}[label=\alph*.]
    \item Écrire les équations des réactions.
    \item Déterminer, en moles, la quantité d'ions hydroxyde apportés par la soude et la quantité d'ions hydronium présents dans la solution d'acide chlorhydrique avant le mélange.
  \end{enumerate}
  \item Le mélange obtenu est-il acide ou basique ?
  \item Calculer les concentrations des ions présents dans le mélange. En déduire le pH.
  \item Quel volume de solution d'acide chlorhydrique de concentration $C_1$ faut-il ajouter au mélange précédent pour obtenir une solution neutre ?
\end{enumerate}
\noindent On donne : $K_e(\text{eau}) = 10^{-14}$ à $25\,^\circ$C ; masses atomiques en g/mol : Na : 23 ; O : 16 ; H : 1.

% ============================ PHYSIQUE ============================
\matiere{PHYSIQUE}{12 points}

\exo{1}{(Oscillations mécaniques)}
\begin{minipage}{0.68\linewidth}
Une tige homogène $OA$ de longueur $L = 1$ m, de masse $m = 100$ g, peut osciller sans frottement autour d'un axe horizontal $(\Delta)$ passant par son extrémité supérieure $O$. On fixe à l'autre extrémité $A$ de la tige une masse $m_A = \dfrac{3m}{2}$.

Le pendule ainsi constitué est écarté de sa position d'équilibre d'un angle $\theta_0 = 0{,}15$ rad, puis il est abandonné sans vitesse initiale.
\end{minipage}\hfill
\begin{minipage}{0.28\linewidth}
\centering
\begin{tikzpicture}[>=Stealth,scale=1]
  \draw[thick] (0,0)--(0,-3.2);
  \node[above] at (0,0){$O$};
  \node[right] at (0.05,-0.05){$(\Delta)$};
  \draw (-0.12,-0.12)--(0.12,0.12); \draw (-0.12,0.12)--(0.12,-0.12);
  \fill (0,-3.2) circle (2pt) node[below]{$A$};
\end{tikzpicture}
\end{minipage}

\begin{enumerate}
  \item Soit $G$ le centre d'inertie du système ainsi constitué.
  \begin{enumerate}[label=\alph*.]
    \item Montrer que la position du centre d'inertie $G$ est telle que $OG = \dfrac{4L}{5}$.
    \item Calculer le moment d'inertie $J_\Delta$ du système par rapport à $(\Delta)$.
  \end{enumerate}
  \item \begin{enumerate}[label=\alph*.]
    \item En utilisant la méthode énergétique, déterminer la nature du mouvement de ce pendule pour des oscillations de faible amplitude.
    \item Écrire l'équation horaire du mouvement de ce pendule en prenant pour origine des temps l'instant où on l'abandonne.
  \end{enumerate}
\end{enumerate}

\exo{2}{(Onde progressive)}
Une lame vibrante est munie d'une pointe dont l'extrémité frappe la surface d'une nappe d'eau en un point $S$. La pointe a un mouvement rectiligne sinusoïdal de fréquence $N = 50$ Hz et d'amplitude $a = 3$ mm.

\begin{enumerate}
  \item Établir l'équation horaire du mouvement de $S$, sachant qu'à $t = 0$, la source $S$ est à son élongation maximale positive.
  \item La nappe d'eau est le siège d'une onde progressive sinusoïdale transversale de longueur d'onde $\lambda = 2$ cm.
  \begin{enumerate}[label=\alph*.]
    \item Calculer la célérité des ondes.
    \item Établir l'équation horaire $Y_M(t)$ du mouvement d'un point $M$ situé à la distance $x = 8{,}5$ cm de $S$.
    \item Comparer les mouvements de $S$ et de $M$.
    \item Représenter dans un même système d'axes les courbes $Y_S(t)$ et $Y_M(t)$.
  \end{enumerate}
\end{enumerate}

\exo{3}{(Effet photoélectrique)}
On éclaire une cellule photoélectrique au césium successivement avec deux radiations lumineuses de longueur d'onde $\lambda_1 = 410$ nm et $\lambda_2 = 740$ nm. On rappelle que $1$ nm $= 10^{-9}$ m.

\begin{enumerate}
  \item La longueur d'onde seuil photoélectrique du césium est $\lambda_0 = 660$ nm.
  \begin{enumerate}[label=\alph*.]
    \item Donner la définition de la longueur d'onde seuil.
    \item Parmi les deux radiations, préciser celle qui provoque l'émission d'électrons.
  \end{enumerate}
  \item Pour la radiation qui provoque l'émission d'électrons, calculer en électron-volts l'énergie cinétique maximale d'un électron émis par la cathode.
  \item On établit entre l'anode et la cathode une tension variable $U_{AC}$ et on note l'intensité du courant pour chaque valeur de $U_{AC}$. On donne la caractéristique $I = f(U_{AC})$.

  \begin{center}
  \begin{tikzpicture}[>=Stealth,scale=1]
    \draw[->] (-2,0)--(6,0) node[right]{$U_{AC}$ (V)};
    \draw[->] (0,-0.3)--(0,2.4) node[above]{$i(\mu$A$)$};
    \draw[dashed] (-1.15,1.7)--(5.6,1.7);
    \node[left] at (0,1.7){$1{,}2$};
    \draw (-1.15,0.05)--(-1.15,-0.05) node[below]{$-1{,}15$};
    \node[below right] at (0,0){$0$};
    \draw[thick] plot[smooth,tension=0.7] coordinates {(-1.15,0)(-0.7,0.12)(-0.2,0.55)(0.3,1.25)(0.9,1.6)(1.8,1.7)(5.4,1.7)};
  \end{tikzpicture}
  \end{center}

  \begin{enumerate}[label=\alph*.]
    \item Que signifient les nombres $-1{,}15$ V et $1{,}2$ µA ?
    \item Calculer la tension $U_{AC}$ pour laquelle les électrons arrivent à l'anode à la vitesse $V_A = 2\,000$ km/s.
    \item Lorsqu'on a obtenu le courant de saturation, déterminer le nombre d'électrons émis par la cathode en 16 s.
  \end{enumerate}
\end{enumerate}

\end{document}
